\section{Araştırma Deneyimi}

\begin{twocolentry}{
    \textit{Nis 2024 – Nis 2025}\\
    \small Madrid, İspanya}
    \normalsize \textbf{Araştırma Stajyeri}, \small {IMDEA Software Institute}
    \small  
    \begin{itemize}
        \item İncelenen sistemlerde static analysis kullanarak dormant vulnerabilities \textit{(CVEs)} tespit etmek için bir metodoloji geliştirdi, özellikle Debian-based containers üzerine odaklandı.
        \item Algılanan bir vulnerability'nin kritik bir tehdit oluşturup oluşturmadığını değerlendiren ve bildiren bir tool tasarlayıp geliştirdi.
      \end{itemize}
\end{twocolentry}

\vspace{\betweenentryspace}
\begin{twocolentry}{
    \textit{Nis 2023 – Tem 2023}\\
    \small Lannion, Fransa}
    \normalsize \textbf{Yazılım Geliştirme Stajyeri}, \small {Nokia Bell Labs}
    \small  
    \begin{itemize}
        \item Guncel cloud FPGA alanında multi-tenant client isolation üzerine devam eden araştırmalar için state-of-the-art bir proof of concept \textit{(PoC)} geliştirdi.
        \item k3s ve Docker kullanarak cloud deployment infrastructure kurdu. \textit{PoC}'yi ekip için dokümante etti. \\
        \footnotesize \textbf{Araçlar:} Golang, Kubernetes (K3s), Docker, Embedded Linux, Petalinux, Yocto
    \end{itemize}
\end{twocolentry}

\vspace{\betweenentryspace}

\begin{twocolentry}{
    \textit{Eki 2022 – Tem 2023}\\
    \small İstanbul, Türkiye}
    \normalsize \textbf{Lisans Tezi}, \small {Sabancı Üniversitesi}
    \small  
    \begin{itemize}
        \item Timing attacks gerçekleştirebilen programları tespit etmek için control flow graphs analiz eden ve suspicious execution paths belirleyen bir static analysis yöntemi geliştirdi.
        \item Değerlendirmeler sonucunda zararlı program tespitinde yüksek doğruluk oranı elde etti ve geniş bir yelpazede zararsız ve zararlı örnek veriler üzerinde güçlü sonuçlar aldı.
    \end{itemize}
\end{twocolentry}


\vspace{\betweenentryspace}

\begin{twocolentry}{
    \textit{Bahar 2020}\\
    \small İstanbul, Türkiye}
    \normalsize \textbf{Lisans Araştırmacısı}, \small \href{https://pure.sabanciuniv.edu/project/4041/solving-matching-problems-using-artificial-intelligence-interview-scheduling}{Sabancı Üniversitesi - PURE}

    \small  
    \begin{itemize}
        \item Bipartite graphs ve perturbations for stable matching üzerine odaklanarak computational methods kullanarak matching problems çözdü.
    \end{itemize}
\end{twocolentry}
